\chapter{Podstawy teorii grafów}
\label{ch:tg-podstawy}

Teoria grafów bada układy obiektów i relacji między nimi. Wierzchołki będą
reprezentować obiekty, a krawędzie ich powiązania. Już na poziomie tego
prostego opisu pojawiają się trzy różne perspektywy: kombinatoryczna
(liczenie połączeń i dróg), algebraiczna (macierze) i topologiczna
(realizacja grafu jako przestrzeni jednowymiarowej). Pierwsze dwa rozdziały
poświęcamy części kombinatorycznej, nie zakładając znajomości żadnego
pojęcia z teorii grafów.

\section{Graf prosty, multigraf i realizacja geometryczna}
\label{sec:tg-obiekty}

\begin{definicja}[Graf prosty]
\emph{Skończonym grafem prostym nieskierowanym} nazywamy parę
$G=(V,E)$, gdzie $V$ jest skończonym zbiorem
\emph{wierzchołków} (dopuszczamy także $V=\varnothing$), natomiast
$E\subseteq\binom V2=\{\{u,v\}:u,v\in V,\ u\ne v\}$
jest zbiorem \emph{krawędzi}. Zapis $uv\in E$ oznacza
$\{u,v\}\in E$. Wierzchołki $u,v$ są \emph{sąsiadujące},
gdy $uv\in E$; oba są wówczas \emph{incydentne}
z krawędzią $uv$.
\end{definicja}

Z definicji żadna krawędź nie łączy wierzchołka z nim samym i żadna para
różnych wierzchołków nie jest połączona więcej niż jedną krawędzią.
Ważne jest jednak, by nie utożsamiać pojęcia grafu z jego rysunkiem:
przecięcia linii na rysunku nie są nowymi wierzchołkami, o ile ich
nie zaznaczono. Dwa różne rysunki mogą opisywać dokładnie ten sam graf.

\begin{definicja}[Multigraf z pętlami]
\emph{Skończonym multigrafem nieskierowanym z pętlami}
nazywamy trójkę $G=(V,E,\partial)$, w której $V$ i $E$ są
skończonymi zbiorami, a $\partial:E\to (V\times V)/\mathfrak S_2$
przyporządkowuje każdej krawędzi nieuporządkowaną parę
\emph{z zachowaniem krotności}. Zapis $\partial(e)=[(u,v)]$
oznacza orbitę pary $(u,v)$ przy zamianie współrzędnych;
w szczególności $[(u,u)]$ jest poprawną parą równych końców.
Jeśli $u=v$, krawędź $e$ nazywamy \emph{pętlą}.
Różne elementy $e,f\in E$ mogą mieć te same końce;
wówczas są \emph{krawędziami równoległymi}.
\end{definicja}

W całej pierwszej części, jeśli nie podano inaczej, \emph{graf} oznacza
graf prosty nieskierowany i skończony. Kiedy mówimy o multigrafie,
powiemy to jawnie. Później multigrafy będą potrzebne m.in.
przy diagramach węzłów i grafach dualnych triangulacji.

\begin{center}
\begin{tikzpicture}[scale=.91,line cap=round]
  \begin{scope}[shift={(-3.05,0)}]
    \node[tgBlueVertex] (a) at (0,0) {$a$};
    \node[tgBlueVertex] (b) at (2.0,0) {$b$};
    \node[tgBlueVertex] (c) at (1,1.55) {$c$};
    \draw[tgEdge] (a)--(b)--(c)--(a);
    \node[font=\small,text=tgNavy] at (1,-.70) {graf prosty};
  \end{scope}
  \begin{scope}[shift={(2.2,0)}]
    \node[tgRedVertex] (u) at (0,0) {$u$};
    \node[tgRedVertex] (v) at (2.0,0) {$v$};
    \draw[tgEdge] (u) to[bend left=26] (v);
    \draw[tgAccent] (u) to[bend right=26] (v);
    \draw[tgEdge] (u) .. controls (-1.10,1.15) and (-1.10,-.38) .. (u);
    \node[font=\small,text=tgNavy,align=center] at (1,-.85)
      {multigraf z pętlą};
  \end{scope}
\end{tikzpicture}

{\footnotesize Po lewej $E=\{ab,bc,ca\}$. Po prawej są dwie różne
krawędzie o końcach $u,v$ oraz pętla przy $u$.}
\end{center}

\begin{definicja}[Graf pełny i graf bez krawędzi]
Dla liczby całkowitej $n\ge 1$ przez $K_n$ oznaczamy graf
\emph{pełny} o wierzchołkach $\{1,\ldots,n\}$ i wszystkich
możliwych krawędziach. Graf \emph{bez krawędzi} na $n$
wierzchołkach ma zbiór krawędzi równy $\varnothing$;
można go oznaczyć przez $\overline{K_n}$.
Od grafu bez krawędzi odróżniamy \emph{graf zerowy}
$G_0=(\varnothing,\varnothing)$. Graf zerowy jest dopuszczalny zgodnie z definicją grafu prostego.
Niektóre pojęcia, np. najmniejszy i największy stopień,
będą jednak wymagały niepustego zbioru wierzchołków.
\end{definicja}

\begin{przyklad}[Liczenie krawędzi grafu pełnego]
Każda z $n$ etykiet sąsiaduje w $K_n$ z $n-1$ pozostałymi.
Licząc uporządkowane pary końców, dostajemy $n(n-1)$,
ale każdą krawędź policzyliśmy dwukrotnie. Stąd
$|E(K_n)|=\binom n2=n(n-1)/2$.
Na przykład $K_4$ ma sześć krawędzi. Rysując go na płaszczyźnie
z dwoma przecinającymi się przekątnymi, nie tworzymy z ich
przecięcia piątego wierzchołka.
\end{przyklad}

\begin{definicja}[Izomorfizm]
Niech $G=(V,E)$ i $H=(W,F)$ będą grafami prostymi.
\emph{Izomorfizmem grafów} $G\to H$ nazywamy bijekcję
$f:V\to W$, dla której
$uv\in E\iff f(u)f(v)\in F$ dla wszystkich różnych $u,v\in V$.
Piszemy wtedy $G\cong H$.
\end{definicja}

Izomorfizm zmienia jedynie nazwy wierzchołków. Wyznacza
bijekcję między zbiorami krawędzi, gdyż zachowuje i odzwierciedla
sąsiedztwo. W szczególności zachowuje liczby wierzchołków
i krawędzi. Dalsze niezmienniki izomorfizmu poznamy po
wprowadzeniu stopni, cykli i spójności.

\begin{definicja}[Realizacja geometryczna]
Dla grafu prostego $G=(V,E)$ rozważmy przestrzeń
$V\sqcup\bigsqcup_{e\in E}[0,1]_e$ i utożsammy
końce odcinka $[0,1]_e$ z dwoma wierzchołkami
należącymi do krawędzi $e$. Otrzymany iloraz
$|G|$ nazywamy \emph{realizacją geometryczną} grafu.
\end{definicja}

Realizacja jest kompleksem CW wymiaru co najwyżej jeden
(jej własności topologiczne omówimy w dalszej części). Jej definicja nie
zależy od tego, jak rysujemy graf na kartce. W szczególności
samo przecinanie się krzywych na rysunku nie prowadzi do
utożsamienia odpowiednich punktów w $|G|$.

\section{Stopnie wierzchołków i podstawowe rachunki}
\label{sec:tg-stopnie}

\begin{definicja}[Stopień]
Dla grafu prostego $G=(V,E)$ i $v\in V$ przez
$N_G(v)=\{u\in V:uv\in E\}$ oznaczamy
\emph{sąsiedztwo} $v$, a liczbę $d_G(v)=|N_G(v)|$
nazywamy \emph{stopniem} $v$. W multigrafie stopień
jest liczbą końców krawędzi incydentnych z danym
wierzchołkiem; każda pętla wnosi do stopnia dwa.
\end{definicja}

W grafie prostym zachodzi $0\le d_G(v)\le |V|-1$.
Stopień zero oznacza wierzchołek izolowany, a stopień jeden
--- liść. Dla skończonego niepustego grafu definiujemy
$\delta(G)=\min_{v\in V}d_G(v)$ i
$\Delta(G)=\max_{v\in V}d_G(v)$.

\begin{twierdzenie}[Lemat o uściskach dłoni]
Niech $G=(V,E,\partial)$ będzie skończonym multigrafem
nieskierowanym, z dopuszczeniem pętli, a $d_G(v)$
niech liczy każdą pętlę dwukrotnie. Wówczas
\[
\sum_{v\in V}d_G(v)=2|E|.
\]
W szczególności liczba wierzchołków nieparzystego
stopnia jest parzysta, także gdy multigraf ma pętle.
\end{twierdzenie}

\begin{proof}
Rozważmy zbiór incydencji, którego elementami są pary
$(v,h)$, gdzie $h$ jest końcem pewnej krawędzi przylegającym
do $v$. Dla każdej krawędzi są dokładnie dwa takie końce,
również wtedy, gdy jest ona pętlą. Liczba wszystkich
incydencji wynosi zatem $2|E|$. Jeśli natomiast
najpierw pogrupujemy incydencje według wierzchołków,
to przy $v$ otrzymamy dokładnie $d_G(v)$ końców.
Oba sposoby zliczania dają pierwszą równość.

Niech $V_{\mathrm{odd}}$ będzie zbiorem wierzchołków
nieparzystego stopnia. Redukując uzyskaną równość modulo
$2$, otrzymujemy
$0\equiv\sum_{v\in V}d_G(v)\equiv
|V_{\mathrm{odd}}|\pmod2$, co dowodzi drugiej części tezy.
\end{proof}

\begin{przyklad}[Obliczenie stopni]
Weźmy graf o wierzchołkach $a,b,c,d,e$ i krawędziach
$ab,ac,ad,bc,cd$. Wówczas
$(d(a),d(b),d(c),d(d),d(e))=(3,2,3,2,0)$.
Suma stopni wynosi $10=2\cdot5$. Dokładnie dwa
wierzchołki mają nieparzysty stopień. Wierzchołek $e$
jest izolowany.
\end{przyklad}

\begin{center}
\begin{tikzpicture}[scale=.96]
  \node[tgBlueVertex] (a) at (0,1.4) {$a$};
  \node[tgRedVertex] (b) at (-1.8,0) {$b$};
  \node[tgRedVertex] (c) at (1.8,0) {$c$};
  \node[tgBlueVertex] (d) at (0,-1.25) {$d$};
  \node[tgTealVertex] (e) at (3.5,1.05) {$e$};
  \draw[tgEdge] (a)--(b)--(c)--(a);
  \draw[tgEdge] (a)--(d)--(c);
  \node[font=\footnotesize,text=tgNavy] at (3.5,.45) {$d(e)=0$};
\end{tikzpicture}

{\footnotesize Przykład z pięcioma krawędziami; izolowany
wierzchołek nie wpływa na sumę stopni.}
\end{center}

\begin{wniosek}
Nie istnieje graf z dokładnie jednym wierzchołkiem
nieparzystego stopnia. Jeżeli graf prosty o $n$
wierzchołkach jest $r$-regularny (czyli każdy wierzchołek
ma stopień $r$), to $nr$ jest parzyste i
$|E|=nr/2$.
\end{wniosek}

\begin{proof}
Obie własności wynikają bezpośrednio z poprzedniego
twierdzenia: odpowiednio z parzystości liczby
wierzchołków nieparzystego stopnia i ze wzoru
$\sum_{v\in V}d(v)=nr=2|E|$.
\end{proof}

\section{Podgrafy, dopełnienie i macierz sąsiedztwa}
\label{sec:tg-podgrafy}

\begin{definicja}[Podgraf i podgraf indukowany]
Niech $G=(V,E)$ będzie grafem prostym. Graf
$H=(W,F)$ jest jego \emph{podgrafem}, jeśli
$W\subseteq V$ i $F\subseteq E\cap\binom W2$.
Dla dowolnego $W\subseteq V$ graf
$G[W]=(W,E\cap\binom W2)$ nazywamy
\emph{podgrafem indukowanym przez $W$}.
\end{definicja}

W podgrafie można usuwać zarówno wierzchołki, jak i część
krawędzi pomiędzy pozostałymi wierzchołkami. W podgrafie
indukowanym zachowujemy \emph{wszystkie} krawędzie między
wybranymi wierzchołkami. Notacja $G-e$ oznacza usunięcie
samej krawędzi $e$; notacja $G-v$ usuwa wierzchołek $v$
wraz z wszystkimi jego krawędziami. W szczególności
$G[\varnothing]=G_0$.

\begin{przyklad}[Podgraf a podgraf indukowany]
Niech $G=K_4$ na zbiorze $V=\{1,2,3,4\}$ i niech
$W=\{1,2,3\}$. Graf $H=(W,\{12,23\})$ jest podgrafem $G$,
ale nie jest podgrafem indukowanym przez $W$, gdyż pomija
krawędź $13$ występującą w $G$. Natomiast
$G[W]=(W,\{12,23,13\})$. Podgraf indukowany przez dany
zbiór wierzchołków jest wyznaczony jednoznacznie.
\end{przyklad}

\begin{center}
\begin{tikzpicture}[scale=.9]
  \begin{scope}[shift={(-2.2,0)}]
    \node[tgBlueVertex] (a) at (-1,0) {$1$};
    \node[tgBlueVertex] (b) at (0,1.4) {$2$};
    \node[tgBlueVertex] (c) at (1,0) {$3$};
    \draw[tgEdge] (a)--(b)--(c);
    \node[font=\small] at (0,-.65) {$H$};
  \end{scope}
  \begin{scope}[shift={(2.2,0)}]
    \node[tgBlueVertex] (d) at (-1,0) {$1$};
    \node[tgBlueVertex] (e) at (0,1.4) {$2$};
    \node[tgBlueVertex] (f) at (1,0) {$3$};
    \draw[tgEdge] (d)--(e)--(f);
    \draw[tgAccent] (d)--(f);
    \node[font=\small] at (0,-.65) {$G[W]$};
  \end{scope}
\end{tikzpicture}

{\footnotesize W podgrafie indukowanym nie można pominąć
krawędzi $13$ istniejącej w grafie wyjściowym.}
\end{center}

\begin{definicja}[Dopełnienie]
Dla grafu prostego $G=(V,E)$ jego
\emph{dopełnienie} $\overline G$ jest grafem na
tym samym zbiorze wierzchołków, w którym
$uv\in E(\overline G)$ wtedy i tylko wtedy,
gdy $u\ne v$ oraz $uv\notin E(G)$.
\end{definicja}

\begin{stwierdzenie}
Dla grafu prostego $G$ o $n$ wierzchołkach oraz
każdego $v\in V(G)$ zachodzi
$d_{\overline G}(v)=n-1-d_G(v)$, a ponadto
$|E(\overline G)|=\binom n2-|E(G)|$.
\end{stwierdzenie}

\begin{proof}
Spośród pozostałych $n-1$ wierzchołków dokładnie
$d_G(v)$ jest sąsiadami $v$ w $G$. Pozostałe
$n-1-d_G(v)$ są i tylko one sąsiadami $v$ w
$\overline G$. Każda nieuporządkowana para różnych
wierzchołków należy dokładnie do jednego ze zbiorów
$E(G)$, $E(\overline G)$, zatem sumy ich mocy dają
$\binom n2$.
\end{proof}

\begin{definicja}[Macierz sąsiedztwa]
Niech $G=(V,E)$ będzie grafem prostym o
$V=\{v_1,\ldots,v_n\}$. \emph{Macierzą sąsiedztwa}
nazywamy $A_G=(a_{ij})\in M_n(\mathbb Z)$,
gdzie $a_{ij}=1$, jeśli $v_iv_j\in E$, oraz
$a_{ij}=0$ w przeciwnym przypadku.
\end{definicja}

Macierz $A_G$ jest symetryczna i ma zerową przekątną.
Sumy wierszy to stopnie wierzchołków.

\begin{przyklad}[Ścieżka na czterech wierzchołkach]
Niech $V=\{1,2,3,4\}$, a $E=\{12,23,34\}$.
Wtedy
\[
A_G=\begin{pmatrix}
0&1&0&0\\1&0&1&0\\0&1&0&1\\0&0&1&0
\end{pmatrix}.
\]
Mnożenie $A_G(1,1,1,1)^T$ daje $(1,2,2,1)^T$,
a więc wektor stopni. Uogólnienie tej obserwacji
poprowadzi nas później do laplasjanu grafu i
jego widma.
\end{przyklad}

\begin{stwierdzenie}[Zliczanie spacerów macierzami]
Niech $G$ będzie grafem prostym o macierzy sąsiedztwa
$A_G$. Dla liczby całkowitej $k\ge0$ element
$(A_G^k)_{ij}$ jest równy liczbie ciągów
$(v_i=v_{i_0},v_{i_1},\ldots,v_{i_k}=v_j)$,
w których każda kolejna para jest krawędzią grafu.
Dla $k=0$ liczymy ciąg zerowej długości tylko wtedy,
gdy $i=j$.
\end{stwierdzenie}

\begin{proof}
Dla $k=0$ twierdzenie jest równością $A_G^0=I$.
Dla $k=1$ jest definicją macierzy sąsiedztwa.
Załóżmy, że jest prawdziwe dla $k$. Z definicji
mnożenia macierzy
$(A_G^{k+1})_{ij}=\sum_{r=1}^n(A_G^k)_{ir}a_{rj}$.
Każdy ciąg długości $k+1$ kończący się w $v_j$
ma dokładnie jeden przedostatni wierzchołek $v_r$.
Wybór poprzednich $k$ kroków daje $(A_G^k)_{ir}$
możliwości, a ostatni krok istnieje dokładnie przy
$a_{rj}=1$. Sumowanie po $r$ zlicza wszystkie
ciągi dokładnie raz. Indukcja jest zakończona.
\end{proof}

Ważna subtelność: takie ciągi mogą wielokrotnie odwiedzać
ten sam wierzchołek lub tę samą krawędź. Nie są więc z zasady
drogami prostymi ani cyklami.

\section{Spacery, szlaki, ścieżki i cykle}
\label{sec:tg-drogi}

\begin{definicja}[Spacer, szlak i ścieżka]
Niech $G=(V,E)$ będzie grafem prostym.
\emph{Spacerem} długości $k\ge0$ z $v_0$ do $v_k$
nazywamy ciąg $v_0,v_1,\ldots,v_k$ taki, że
$v_{i-1}v_i\in E$ dla $i=1,\ldots,k$.
Spacer nazywamy \emph{szlakiem}, jeśli żadna
krawędź nie jest przechodzona dwukrotnie.
Nazywamy go \emph{ścieżką} (\emph{drogą prostą}), jeśli
wierzchołki $v_0,\ldots,v_k$ są parami różne.
\emph{Spacer zamknięty} ma $v_0=v_k$.
\end{definicja}

Gdy w dalszym tekście pojawi się słowo ,,droga'' w znaczeniu
kombinatorycznym, będziemy mieli na myśli drogę prostą. W innych kontekstach matematycznych
,,droga'' może oznaczać dowolne ciągłe odwzorowanie odcinka;
rozróżnienie to jest istotne.

\begin{definicja}[Cykl]
\emph{Cyklem} długości $k\ge3$ w grafie prostym
nazywamy spacer zamknięty $v_0,v_1,\ldots,v_k=v_0$,
w którym $v_0,\ldots,v_{k-1}$ są parami różne.
Dwa zapisy powstałe przez zmianę punktu startowego
lub kierunku obiegu wyznaczają ten sam cykl.
Graf \emph{acykliczny} nie zawiera cyklu.
\end{definicja}

Dla $n\ge1$ oznaczamy przez $P_n$ graf o wierzchołkach
$1,\ldots,n$ i krawędziach $\{i,i+1\}$ dla $1\le i<n$.
Dla $n\ge3$ graf $C_n$ powstaje z $P_n$ przez dodanie
krawędzi $\{1,n\}$.

\begin{stwierdzenie}[Zliczanie trójkątów]
\label{prop:tg-trace-triangles}
Niech $G$ będzie skończonym grafem prostym o co najmniej
jednym wierzchołku i macierzy sąsiedztwa $A_G$.
Liczba jego cykli długości trzy wynosi
$\operatorname{tr}(A_G^3)/6$, gdzie $\operatorname{tr}$
oznacza sumę elementów diagonalnych macierzy.
\end{stwierdzenie}

\begin{proof}
Ze wzoru zliczającego spacery wynika, że
$\operatorname{tr}(A_G^3)$ zlicza spacery zamknięte
długości trzy z wyróżnionym początkiem. W każdym takim
spacerze $u,v,w,u$ wierzchołki $u,v,w$ są parami różne,
gdyż każda równość dwóch z nich wymuszałaby pętlę.
Każdy taki spacer wyznacza trójkąt, a każdy trójkąt
daje dokładnie sześć spacerów: trzy wybory początku
i dwa kierunki obiegu. Stąd podany wzór.
\end{proof}

\begin{przyklad}
Dla $K_3$ każdy element diagonalny $A_G^3$ wynosi dwa,
więc ślad wynosi sześć, a liczba trójkątów jeden.
Dla $P_4$ ślad wynosi zero.
\end{przyklad}


W multigrafie spacer trzeba opisać naprzemiennym ciągiem
$v_0,e_1,v_1,\ldots,e_k,v_k$, gdzie
$\partial(e_i)=[(v_{i-1},v_i)]$. Sam ciąg wierzchołków
nie rozróżnia krawędzi równoległych. Szlak nie powtarza
krawędzi; cykl to zamknięty szlak dodatniej długości,
w którym $v_0,\ldots,v_{k-1}$ są parami różne.
Dopuszcza to cykle długości jeden (pętle) i dwa
(dwie różne krawędzie równoległe).

\begin{lemat}[Usuwanie zamkniętych fragmentów spaceru]
Niech $G$ będzie grafem skończonym i niech
$v_0,\ldots,v_k$ będzie spacerem z $u=v_0$
do $v=v_k$, gdzie $u\ne v$. Wtedy w $G$
istnieje droga prosta z $u$ do $v$ długości
nie większej niż $k$.
\end{lemat}

\begin{proof}
Jeśli w spacerze żaden wierzchołek się nie powtarza,
to jest on już drogą prostą. W przeciwnym razie istnieją
$i<j$ takie, że $v_i=v_j$. Usuwamy fragment od
$v_{i+1}$ do $v_j$: otrzymujemy spacer
$v_0,\ldots,v_i,v_{j+1},\ldots,v_k$
(jeśli $j=k$, spacer kończy się w $v_i=v$).
Zachowuje on oba końce i jest krótszy o $j-i$ kroków.
Powtarzamy tę operację. Długość jest nieujemną liczbą
całkowitą i za każdym razem maleje, więc proces
musi się zakończyć. Pozostały spacer nie powtarza
wierzchołków, czyli jest drogą prostą.
\end{proof}

\begin{uwaga}
Nie wolno w powyższym lemacie zastąpić ,,spaceru między
dwoma różnymi końcami'' przez dowolny niepusty spacer
zamknięty i wnioskować, że istnieje cykl. W grafie złożonym z jednej krawędzi $uv$
spacer $u,v,u$ jest zamknięty, ale nie zawiera cyklu.
Ten sam fragment krawędzi został przebyty dwa razy.
\end{uwaga}

\section{Spójność i składowe spójne}
\label{sec:tg-spojnosc}

\begin{definicja}[Spójność]
Wierzchołki $u,v$ grafu $G$ nazywamy
\emph{połączonymi}, jeśli $u=v$ albo istnieje
spacer od $u$ do $v$. Graf $G$ jest \emph{spójny},
jeśli $V\ne\varnothing$ oraz każda para jego wierzchołków jest połączona.
Maksymalne względem inkluzji niepuste zbiory parami połączonych
wierzchołków, wraz z wszystkimi krawędziami między nimi,
nazywamy \emph{składowymi spójnymi} grafu.
\end{definicja}

\begin{stwierdzenie}
Dla każdego grafu $G=(V,E)$ relacja
$u\sim v\iff u,v$ są połączone jest relacją
równoważności na $V$. Jej klasami są dokładnie
zbiory wierzchołków składowych spójnych.
\end{stwierdzenie}

\begin{proof}
Relacja jest zwrotna, bo dopuszczamy spacer długości zero.
Jest symetryczna, bo każdy spacer można przejść w kierunku
odwrotnym. Jeśli spacer łączy $u$ z $v$, a drugi łączy
$v$ z $w$, ich sklejenie daje spacer z $u$ do $w$,
więc relacja jest przechodnia. Klasy relacji są
parami rozłączne i pokrywają $V$. Jeżeli wierzchołek
spoza danej klasy miałby krawędź do wierzchołka tej klasy,
byłby z nim połączony i sam należałby do klasy. Każdy spacer zaczynający się w danej klasie pozostaje
w niej, ponieważ każdy odwiedzony wierzchołek jest połączony
z początkiem spaceru. Dlatego graf indukowany przez klasę
jest spójny. Każdy niepusty zbiór parami połączonych
wierzchołków zawiera się w jednej klasie, a więc maksymalne
takie zbiory są dokładnie klasami równoważności.
W szczególności graf zerowy ma zero składowych.
\end{proof}

\begin{przyklad}[Spójny podgraf nie musi być składową]
W cyklu $C_5$ dwa sąsiadujące wierzchołki wraz z
łączącą je krawędzią tworzą spójny podgraf. Nie jest
on jednak składową spójną całego $C_5$, ponieważ
pozostałe trzy wierzchołki można do niego dołączyć,
nie tracąc spójności. Składowa spójna oznacza
\emph{maksymalny} spójny fragment.
\end{przyklad}

\begin{center}
\begin{tikzpicture}[scale=.91]
  \begin{scope}[shift={(-3.0,0)}]
    \node[tgBlueVertex] (a) at (-.9,-.4) {$a$};
    \node[tgBlueVertex] (b) at (1.0,-.4) {$b$};
    \node[tgBlueVertex] (c) at (0,1.0) {$c$};
    \draw[tgEdge] (a)--(b)--(c)--(a);
    \draw[draw=tgBlue!50,rounded corners=10pt,line width=.9pt]
      (-1.48,-.92) rectangle (1.58,1.55);
  \end{scope}
  \begin{scope}[shift={(1.9,0)}]
    \node[tgRedVertex] (d) at (0,0) {$d$};
    \node[tgRedVertex] (e) at (1.65,0) {$e$};
    \draw[tgAccent] (d)--(e);
  \end{scope}
  \node[tgTealVertex] (f) at (4.65,1.2) {$f$};
  \node[font=\footnotesize,text=tgNavy] at (0,-1.5) {trzy składowe spójne};
\end{tikzpicture}
\end{center}

\begin{definicja}[Odległość]
Dla grafu $G$ i $u,v\in V(G)$ definiujemy
$d_G(u,v)$ jako najmniejszą długość drogi
prostej z $u$ do $v$, jeżeli taka droga istnieje,
oraz $d_G(u,v)=\infty$ w przeciwnym przypadku.
Przyjmujemy $d_G(v,v)=0$.
\end{definicja}

\begin{stwierdzenie}
Na zbiorze wierzchołków każdej składowej spójnej
grafu $G$ funkcja $d_G$ jest metryką.
\end{stwierdzenie}

\begin{proof}
Odległości są nieujemne, a odległość zero zachodzi
tylko dla identycznych wierzchołków. Odwrócenie
drogi daje symetrię. Dla $u,v,w$ w tej samej
składowej wybieramy najkrótsze drogi od $u$ do $v$
i od $v$ do $w$. Ich konkatenacja jest spacerem od
$u$ do $w$ długości $d_G(u,v)+d_G(v,w)$.
Jeśli $u=w$, nierówność trójkąta wynika z nieujemności
odległości. Jeśli $u\ne w$, na mocy lematu o usuwaniu
powtórzeń istnieje droga prosta nie dłuższa niż ten spacer. Zatem
$d_G(u,w)\le d_G(u,v)+d_G(v,w)$.
\end{proof}

W ten sposób nawet skończony graf otrzymuje naturalną
geometrię metryczną. W późniejszej części porównamy tę
metrykę z metryką długości na realizacji $|G|$.

\section{Grafy dwudzielne i parzystość cykli}
\label{sec:tg-dwudzielne}

\begin{definicja}[Graf dwudzielny]
Graf $G=(V,E)$ jest \emph{dwudzielny}, jeśli istnieją
rozłączne zbiory $V_0,V_1$ o sumie $V$ takie, że
każda krawędź ma jeden koniec w $V_0$, a drugi
w $V_1$. Parę $(V_0,V_1)$ nazywamy \emph{podziałem
dwudzielnym}. \emph{Pełny graf dwudzielny} $K_{r,s}$,
gdzie $r,s\ge1$, łączy wszystkimi możliwymi krawędziami
$r$ wierzchołków jednej klasy z $s$ wierzchołkami
drugiej klasy i nie zawiera innych krawędzi.
\end{definicja}

\begin{lemat}[Nieparzysty spacer zamknięty]
\label{lem:tg-odd-closed-walk}
Niech $G$ będzie grafem prostym, a $W$ spacerem zamkniętym
nieparzystej długości w $G$. Podgraf złożony z wierzchołków
i krawędzi występujących w $W$ zawiera cykl nieparzystej długości.
\end{lemat}

\begin{proof}
W tym podgrafie wybierzmy nieparzysty spacer zamknięty
$x_0,\ldots,x_m=x_0$ o najmniejszej długości. Taki spacer
istnieje, ponieważ sam $W$ spełnia wymagania.
Jeżeli $x_i=x_j$ dla $0\le i<j<m$, rozdzielamy spacer
na $x_i,\ldots,x_j$ oraz
$x_0,\ldots,x_i,x_{j+1},\ldots,x_m$.
Oba są zamknięte i mają dodatnie długości $j-i$
oraz $m-(j-i)$. Dokładnie jedna z nich jest nieparzysta,
a obie są mniejsze od $m$, co przeczy minimalności.
Wierzchołki $x_0,\ldots,x_{m-1}$ są więc parami różne.
Brak pętli wyklucza $m=1$, zatem $m\ge3$ i wybrany
spacer jest cyklem nieparzystej długości.
\end{proof}

\begin{twierdzenie}[Kryterium cykli nieparzystych]
Dla skończonego grafu prostego $G$ następujące
warunki są równoważne:
\begin{enumerate}[label=(\roman*)]
\item $G$ jest dwudzielny;
\item $G$ nie zawiera cyklu nieparzystej długości.
\end{enumerate}
\end{twierdzenie}

\begin{proof}
Załóżmy (i). Idąc po kolejnych krawędziach, przy
każdym kroku zmieniamy klasę podziału z $V_0$ na
$V_1$ lub odwrotnie. Powrót do punktu początkowego
może więc nastąpić jedynie po parzystej liczbie
kroków. Żaden cykl nie ma długości nieparzystej.

Załóżmy teraz (ii). Dla grafu zerowego przyjmujemy
$V_0=V_1=\varnothing$. W pozostałych przypadkach
wystarczy skonstruować podział każdej składowej osobno:
sumy odpowiednich klas dają podział całego grafu,
ponieważ nie ma krawędzi pomiędzy składowymi.
W ustalonej składowej wybierzmy wierzchołek $r$
i dla każdego jej wierzchołka $v$ oznaczmy przez
$\ell(v)=d_G(r,v)$ odległość od $r$. Ustalmy
$V_0=\{v:\ell(v)\text{ jest parzyste}\}$
oraz $V_1=\{v:\ell(v)\text{ jest nieparzyste}\}$.
Przypuśćmy, że istnieje krawędź $uv$ z końcami
o tej samej parzystości odległości. Wybierzmy najkrótsze drogi od $r$ do $u$ i od $r$ do $v$.
Przejście pierwszą drogą, następnie krawędzią $uv$, a wreszcie
drugą drogą w kierunku odwrotnym daje spacer zamknięty
długości $\ell(u)+1+\ell(v)$, która jest nieparzysta.

Na mocy lematu~\ref{lem:tg-odd-closed-walk} graf zawiera
cykl nieparzystej długości, sprzecznie z (ii).
Zatem każda krawędź łączy wierzchołki o
różnych parzystościach, co dowodzi (i).
\end{proof}

\begin{center}
\begin{tikzpicture}[scale=.83]
 \begin{scope}[shift={(-2.8,0)}]
  \node[tgBlueVertex] (a) at (-.9,-.75) {$a$};
  \node[tgRedVertex] (b) at (.9,-.75) {$b$};
  \node[tgRedVertex] (c) at (-.9,.9) {$c$};
  \node[tgBlueVertex] (d) at (.9,.9) {$d$};
  \draw[tgEdge] (a)--(b)--(d)--(c)--(a);
  \node[font=\small,text=tgNavy] at (0,-1.35) {cykl $C_4$: dwudzielny};
 \end{scope}
 \begin{scope}[shift={(2.7,0)}]
  \node[tgBlueVertex] (u) at (0,1) {$u$};
  \node[tgRedVertex] (v) at (-1.12,-.8) {$v$};
  \node[tgRedVertex] (w) at (1.12,-.8) {$w$};
  \draw[tgEdge] (u)--(v)--(w)--(u);
  \node[font=\small,text=tgNavy] at (0,-1.35) {$C_3$: brak podziału};
 \end{scope}
\end{tikzpicture}
\end{center}

\begin{przyklad}[Szybkie rozpoznawanie dwudzielności]
Każdy graf acykliczny jest dwudzielny. W szczególności
dotyczy to drzew, czyli spójnych grafów acyklicznych,
które omówimy w następnym rozdziale. Pełny graf $K_{r,s}$ jest dwudzielny bezpośrednio
z definicji. Dla $n\ge3$ cykl $C_n$ jest dwudzielny
dokładnie wtedy, gdy $n$ jest parzyste. W szczególności
$K_3$ i $K_5$ nie są dwudzielne.
\end{przyklad}

\begin{przyklad}[Niezgodność wykryta przez trójkąty]
Graf graniastosłupa trójkątnego (dwa rozłączne trójkąty
połączone trzema krawędziami między odpowiadającymi sobie
wierzchołkami) i graf pełny dwudzielny $K_{3,3}$ mają po sześć wierzchołków, dziewięć krawędzi i wszystkie stopnie
równe trzy. Pierwszy zawiera cykl długości trzy, natomiast drugi
nie zawiera żadnego cyklu nieparzystego. Nie są więc izomorficzne.
Izomorfizm zachowuje stopnie (indukuje bijekcje sąsiedztw)
i cykle (przenosi ciągi różnych wierzchołków z zachowaniem
krawędzi). To przykład, dlaczego sam ciąg stopni nie klasyfikuje grafów.
\end{przyklad}

\begin{center}
\begin{tikzpicture}[scale=.83]
  \begin{scope}[shift={(-2.7,0)}]
    \node[tgBlueVertex] (a) at (-1.7,-1.2) {};
    \node[tgBlueVertex] (b) at (1.7,-1.2) {};
    \node[tgBlueVertex] (c) at (0,1.8) {};
    \node[tgRedVertex] (d) at (-.65,-.43) {};
    \node[tgRedVertex] (e) at (.65,-.43) {};
    \node[tgRedVertex] (f) at (0,.7) {};
    \draw[tgEdge] (a)--(b)--(c)--(a) (d)--(e)--(f)--(d);
    \draw[tgAccent] (a)--(d) (b)--(e) (c)--(f);
    \node[font=\small] at (0,-1.95) {graniastosłup trójkątny};
  \end{scope}
  \begin{scope}[shift={(2.7,0)}]
    \foreach \i in {1,2,3}{
      \node[tgBlueVertex] (l\i) at (-1.1,{1.2*(2-\i)}) {};
      \node[tgRedVertex] (r\i) at (1.1,{1.2*(2-\i)}) {};
    }
    \foreach \i in {1,2,3}{
      \foreach \j in {1,2,3}{\draw[tgEdge] (l\i)--(r\j);}
    }
    \node[font=\small] at (0,-1.95) {$K_{3,3}$};
  \end{scope}
\end{tikzpicture}

{\footnotesize Te same stopnie, różne liczby trójkątów.
Przecięcia krawędzi na prawym rysunku nie są wierzchołkami.}
\end{center}

\section{Które ciągi liczb są ciągami stopni?}
\label{sec:tg-havel}

Nie każdy $n$-elementowy ciąg liczb całkowitych z przedziału
$\{0,\ldots,n-1\}$ jest ciągiem stopni grafu prostego.
Warunek parzystości sumy jest konieczny, lecz nie wystarcza.
Na przykład $(3,3,1,1)$ ma parzystą sumę $8$, ale nie jest
ciągiem stopni grafu prostego na czterech wierzchołkach:
dwa wierzchołki stopnia trzy muszą być połączone ze wszystkimi
innymi, więc pozostałe miałyby stopnie co najmniej dwa.

\begin{definicja}[Ciąg grafowy]
Skończony ciąg nieujemnych liczb całkowitych
$(d_1,\ldots,d_n)$ nazywamy \emph{grafowym},
jeśli istnieje graf prosty o wierzchołkach
$v_1,\ldots,v_n$, dla którego $d_G(v_i)=d_i$
dla każdego $i$.
\end{definicja}

\begin{twierdzenie}[Kryterium Havla--Hakimiego]
Niech $n\ge2$ i niech
$d_1\ge d_2\ge\cdots\ge d_n\ge0$ będą liczbami
całkowitymi. Jeśli $d_1>n-1$, ciąg nie jest grafowy.
Jeżeli $0\le d_1\le n-1$, to ciąg
$(d_1,\ldots,d_n)$ jest grafowy wtedy i tylko
wtedy, gdy ciąg
\[
(d_2-1,\ldots,d_{d_1+1}-1,
  d_{d_1+2},\ldots,d_n)
\]
jest grafowy (po ewentualnym przestawieniu elementów).
Jeśli w ciągu zredukowanym pojawi się liczba ujemna,
uznajemy, że nie jest grafowy.
Jeśli $d_1=0$, wszystkie wyrazy są zerami i ciąg
jest grafowy. Przy stosowaniu redukcji kończymy również
na ciągu pustym (realizowanym przez graf zerowy) lub
jednoelementowym: tylko $(0)$ jest wtedy grafowy.
\end{twierdzenie}

\begin{proof}
Jeśli $d_1>n-1$, wierzchołek prostego grafu
nie może mieć wymaganego stopnia. Załóżmy więc
$0\le d_1\le n-1$. Oznaczmy $d=d_1$.
W jedną stronę, jeśli istnieje graf na wierzchołkach
$v_2,\ldots,v_n$ o stopniach odpowiednio
$d_2-1,\ldots,d_{d+1}-1,d_{d+2},\ldots,d_n$,
to dołączamy nowy wierzchołek $v_1$ i krawędzie
$v_1v_2,\ldots,v_1v_{d+1}$. Stopnie wszystkich
wierzchołków stają się równe zadanym liczbom.

W trudniejszą stronę załóżmy, że istnieje graf $G$
o stopniach $d_1,\ldots,d_n$. Wykażemy za pomocą
\emph{zamiany dwóch krawędzi} (ang.\ $2$-switch), że
można go zmodyfikować bez zmiany stopni tak,
aby $v_1$ był sąsiadem dokładnie $v_2,\ldots,v_{d+1}$.
Wśród wszystkich realizacji z zadanymi stopniami
wybierzmy taką, w której $v_1$ ma najwięcej
sąsiadów w tym wyróżnionym zbiorze.

Gdyby nie był sąsiadem pewnego $v_j$ z indeksu
$2\le j\le d+1$, musiałby być sąsiadem pewnego
$v_k$ o $k>d+1$. Ponieważ $d_j\ge d_k$, pokażemy,
że da się zamienić te sąsiedztwa. Istnieje
wierzchołek $z\notin\{v_1,v_j,v_k\}$ taki,
że $v_jz\in E(G)$, lecz $v_kz\notin E(G)$.
Istotnie, rozważmy sąsiedztwa z pominięciem
wzajemnej krawędzi $v_jv_k$, jeżeli ona występuje.
Gdyby wszyscy pozostali sąsiedzi $v_j$ byli też
sąsiadami $v_k$, to $v_k$ miałby co najmniej
tylu takich sąsiadów co $v_j$ oraz dodatkowo
$v_1$ (który nie sąsiaduje z $v_j$).
Krawędź $v_jv_k$, jeśli jest obecna, zwiększa oba
stopnie po jednej jednostce. Dostajemy więc
$d_k\ge d_j+1$, sprzecznie z założeniem
$d_j\ge d_k$. Zatem szukany $z$ istnieje.

Usuwamy krawędzie $v_1v_k$ i $v_jz$ oraz
dodajemy $v_1v_j$ i $v_kz$. Nowych krawędzi
wcześniej nie było, więc nadal mamy graf prosty.
Każdy z czterech wierzchołków traci jedną i zyskuje
jedną krawędź, zatem wszystkie stopnie pozostają
bez zmian. Liczba sąsiadów $v_1$ w zbiorze
$\{v_2,\ldots,v_{d+1}\}$ wzrosła o jeden,
sprzecznie z wyborem realizacji. Stąd istnieje
realizacja, w której $v_1$ jest połączony
dokładnie z tymi $d$ wierzchołkami.
Usunięcie $v_1$ daje graf o stopniach wskazanych
w twierdzeniu. To kończy dowód.
\end{proof}

\begin{center}
\begin{minipage}{\linewidth}
\centering
\begin{tikzpicture}[scale=.88]
  \begin{scope}[shift={(-2.55,0)}]
    \node[tgBlueVertex] (a) at (-1,.65) {$v_1$};
    \node[tgBlueVertex] (b) at (1,.65) {$v_k$};
    \node[tgBlueVertex] (c) at (-1,-.65) {$v_j$};
    \node[tgBlueVertex] (d) at (1,-.65) {$z$};
    \draw[tgEdge] (a)--(b) (c)--(d);
    \node[font=\small] at (0,-1.25) {przed zamianą};
  \end{scope}
  \draw[->,thick] (-.65,0)--(.65,0);
  \begin{scope}[shift={(2.55,0)}]
    \node[tgBlueVertex] (e) at (-1,.65) {$v_1$};
    \node[tgBlueVertex] (f) at (1,.65) {$v_k$};
    \node[tgBlueVertex] (g) at (-1,-.65) {$v_j$};
    \node[tgBlueVertex] (h) at (1,-.65) {$z$};
    \draw[tgAccent] (e)--(g) (f)--(h);
    \node[font=\small] at (0,-1.25) {po zamianie};
  \end{scope}
\end{tikzpicture}

{\footnotesize Usuwamy dwie krawędzie poziome i dodajemy
dwie pionowe. Pozostałe krawędzie (niewidoczne na rysunku)
nie zmieniają się. Każdy z czterech stopni pozostaje taki sam.}
\end{minipage}
\end{center}

\begin{przyklad}[Obliczenie metodą Havla--Hakimiego]
Rozważmy ciąg $(3,3,2,2,2)$. Po usunięciu
pierwszej trójki zmniejszamy trzy kolejne wyrazy o jeden:
$(2,1,1,2)$. Sortując, otrzymujemy $(2,2,1,1)$.
Powtarzamy: usunięcie pierwszej dwójki daje
$(1,0,1)$, a po uporządkowaniu $(1,1,0)$.
Kolejny krok prowadzi do $(0,0)$, a zatem
pierwotny ciąg jest grafowy. Zachowując etykiety podczas
sortowania, otrzymujemy także konkretną realizację.
W pierwszym kroku wybieramy krawędzie $12,13,14$,
a pozostałe stopnie wierzchołków $2,5,3,4$ to $2,2,1,1$.
W drugim wybieramy $25,23$, a w trzecim $45$.
Zbiór $\{12,13,14,25,23,45\}$ daje dokładnie
stopnie $(3,3,2,2,2)$.
\end{przyklad}

\subsection{Kryterium Erdősa--Gallaia}
\label{sec:tg-erdos-gallai}

Kryterium Havla--Hakimiego pozwala redukować ciąg i odtwarzać
graf. Można też rozpoznać ciąg grafowy bez kolejnych redukcji,
badając nierówności wynikające z liczby możliwych krawędzi.

\begin{twierdzenie}[Kryterium Erdősa--Gallaia]
\label{thm:tg-erdos-gallai}
Niech $n\ge1$, a $d_1\ge\cdots\ge d_n\ge0$ będą liczbami
całkowitymi. Ciąg $(d_1,\ldots,d_n)$ jest grafowy wtedy i tylko
wtedy, gdy $\sum_{i=1}^n d_i$ jest parzysta oraz dla
każdego $k\in\{1,\ldots,n\}$ zachodzi nierówność
\begin{equation}
\label{eq:tg-erdos-gallai}
\sum_{i=1}^k d_i\le k(k-1)+\sum_{i=k+1}^n\min(k,d_i).
\end{equation}
\end{twierdzenie}

\begin{proof}
Jeżeli graf realizuje ciąg, parzystość wynika z lematu
o uściskach dłoni. Niech $T=\{v_1,\ldots,v_k\}$.
Suma stopni na $T$ liczy dwukrotnie krawędzie wewnątrz $T$
i jednokrotnie krawędzie między $T$ a jego dopełnieniem.
Pierwszy wkład nie przekracza $2\binom{k}{2}$, drugi
zaś jest nie większy niż suma $\min(k,d_i)$ po $i>k$.
Daje to~\eqref{eq:tg-erdos-gallai}.

Dla dowodu wystarczalności zastosujemy metodę częściowych
realizacji Tripathiego, Venugopalan i Westa.
Pracujemy na ustalonym zbiorze $v_1,\ldots,v_n$.
Dla grafu $H$ piszemy $a_i=d_H(v_i)$.
Częściową realizacją nazywamy graf spełniający $a_i\le d_i$
dla wszystkich $i$. Indukcyjnie, dla $r=1,\ldots,n$,
wykażemy istnienie częściowej realizacji, w której
$a_i=d_i$ dla $i\le r$, a zbiór wierzchołków o indeksach
większych niż $r$ nie zawiera żadnej krawędzi.

Na początku kroku $r$ istnieje częściowa realizacja
z $a_i=d_i$ dla $i<r$ i bez krawędzi pomiędzy
wierzchołkami ze zbioru $S=\{v_{r+1},\ldots,v_n\}$:
dla $r=1$ wystarcza graf bez krawędzi, a dla większego $r$
korzystamy z poprzedniego kroku. Spośród takich grafów
wybieramy $H$ maksymalizujący $a_r$. Maksimum istnieje,
bo liczba grafów na ustalonym zbiorze jest skończona.
Załóżmy dla sprzeczności, że $a_r<d_r$.

\smallskip
\noindent\textit{Krok A: wszystkie wcześniejsze wierzchołki
są sąsiadami $v_r$.}
Gdyby $i<r$ i $v_iv_r$ nie było krawędzią, to
$a_i=d_i\ge d_r>a_r$. Istnieje więc
$u\in N_H(v_i)\setminus N_H(v_r)$; wierzchołek $u$
jest różny od $v_r$, ponieważ $v_iv_r$ nie jest krawędzią.
Jeśli $d_r-a_r\ge2$, usuwamy $v_iu$ i dodajemy
$v_ru,v_rv_i$. Stopnie poza $v_r$ są zachowane,
a $a_r$ rośnie o dwa bez przekroczenia $d_r$.
Nie powstaje krawędź wewnątrz $S$.

Jeśli $d_r-a_r=1$, parzystość sumy niedoborów
$\sum_i(d_i-a_i)=\sum_i d_i-2|E(H)|$ wymusza
istnienie $j>r$ z $a_j<d_j$. Gdy $v_rv_j$ nie jest
krawędzią, dodajemy ją. W przeciwnym razie
usuwamy $v_iu$ i $v_rv_j$, dodając $v_ru$ i $v_rv_i$.
W tej drugiej operacji $u\ne v_j$, gdyż $v_j$ należy
do $N_H(v_r)$, a $u$ nie. Stopień $v_j$ maleje o jeden,
pozostałe stopnie poza $v_r$ nie zmieniają się,
a $a_r$ wzrasta o jeden. Wszystkie operacje pozostają
w klasie rozważanych grafów i przeczą maksymalności.

\smallskip
\noindent\textit{Krok B: dla $j>r$ zachodzi
$a_j=\min(r,d_j)$.}
Brak krawędzi wewnątrz $S$ daje $a_j\le r$.
Gdyby $a_j<\min(r,d_j)$, to $a_j<d_j$ i $a_j<r$.
Jeżeli $v_rv_j$ nie jest krawędzią, dodajemy ją.
Jeżeli jest krawędzią, pewien $i<r$ nie sąsiaduje
z $v_j$. Ponieważ $a_i>a_r$, istnieje
$u\in N_H(v_i)\setminus N_H(v_r)$ różny od $v_r$:
istotnie, zbiory $N_H(v_i)\setminus\{v_r\}$
i $N_H(v_r)\setminus\{v_i\}$ mają odpowiednio
$a_i-1$ oraz $a_r-1$ elementów. Pierwszy jest
liczniejszy, więc zawiera element spoza drugiego.
Element ten nie może być równy $v_i$, ponieważ
graf jest prosty, ani $v_r$, który został usunięty.
Wobec tego nie należy także do całego $N_H(v_r)$.
Na mocy kroku A wierzchołek $u$ należy do $S$,
a $u\ne v_j$, ponieważ $v_j$ sąsiaduje z $v_r$.
Usuwamy $v_iu$ i dodajemy $v_ru,v_iv_j$.
Stopnie $v_i,u$ pozostają takie same, stopnie
$v_r,v_j$ wzrastają o jeden, nie przekraczając
zadanych wartości. To znów przeczy maksymalności.

\smallskip
\noindent\textit{Krok C: wierzchołki $v_1,\ldots,v_r$
tworzą graf pełny.}
Na mocy kroku A pozostaje sprawdzić krawędzie
$v_iv_j$ dla $i<j<r$. Jeżeli którejś brakuje,
z nierówności $a_i,a_j>a_r$ i wzajemnego sąsiedztwa
$v_r$ z $v_i,v_j$ znajdujemy
$u\in N_H(v_i)\setminus(N_H(v_r)\cup\{v_r\})$ oraz
$w\in N_H(v_j)\setminus(N_H(v_r)\cup\{v_r\})$.
Istnienie wynika z tego samego porównania liczności
co w kroku B. Oba wierzchołki należą do $S$.
Usuwamy $v_iu,v_jw$ i dodajemy $v_iv_j,v_ru$.
Nawet jeśli $u=w$, usuwane krawędzie są różne;
wtedy stopień $u$ maleje o jeden. Jeśli $u\ne w$,
stopień $u$ jest zachowany, a stopień $w$ maleje
o jeden. Stopnie $v_i,v_j$ są zachowane, a $a_r$
wzrasta o jeden. Ponownie otrzymujemy sprzeczność.

Z kroków B i C wynika teraz dokładna równość
\[
\sum_{i=1}^r a_i=r(r-1)+\sum_{i=r+1}^n\min(r,d_i).
\]
Lewa strona jest ściśle mniejsza od $\sum_{i=1}^r d_i$,
ponieważ wcześniejsze stopnie zostały już zrealizowane,
ale $a_r<d_r$. Przeczy to~\eqref{eq:tg-erdos-gallai}
dla $k=r$. Zatem $a_r=d_r$. Indukcja kończy się dla $r=n$
grafem realizującym cały ciąg.
\end{proof}

\begin{przyklad}
Dla ciągu $(3,3,1,1)$ nierówność dla $k=2$ wymagałaby
$6\le2+1+1=4$, co jest fałszem. Parzystość sumy nie
wystarcza więc do realizowalności. Dla $(3,3,2,2,2)$
lewe strony dla kolejnych $k$ to $3,6,8,10,12$,
a prawe to $4,8,10,14,20$. Wszystkie nierówności są
spełnione i suma jest parzysta, zgodnie z wcześniejszą
konstrukcją Havla--Hakimiego.
\end{przyklad}

\begin{uwaga}
Kryterium nierówności daje bezpośredni test dla
uporządkowanego ciągu, podczas gdy redukcja
Havla--Hakimiego dostarcza również prostego sposobu
konstrukcji grafu. Dowód wystarczalności powyżej także
ma charakter konstrukcyjny. Źródło metody:
A.~Tripathi, S.~Venugopalan, D.~B.~West,
\emph{A short constructive proof of the Erdős--Gallai
characterization of graphic lists}, Discrete Mathematics
310 (2010), 843--844.
\end{uwaga}

\begin{uwaga}[Zakres pierwszego rozdziału]
Dotychczas nie używaliśmy topologii algebraicznej ani
narzędzi różniczkowych. Za to każdy z poznanych obiektów
ma interpretację topologiczną: cykl staje się okręgiem,
składowe spójne mają sens dla realizacji $|G|$, a
właściwości drzew okażą się związane z trywialnością
pętli. W następnym rozdziale przygotujemy te związki
na poziomie czysto kombinatorycznym.
\end{uwaga}

